\chapter{Discrete Distributions}

In the last chapter you met the random variable: a number whose value is decided
by chance, with a probability mass function (PMF) that tells you how much weight
sits on each possible value. So far each PMF has been a one-off, invented for the
problem at hand. But a small handful of PMFs show up again and again, across coin
games, genetics, call centers, quality control, web traffic, radioactivity, and
machine learning, and they are not arbitrary. Each one is the answer to a
\emph{story} -- a particular way of generating randomness -- and once you know the
story, you can recognize the distribution when it appears and write down its mean,
its variance, and its shape without re-deriving anything.

We will meet six named families: \textbf{Bernoulli}, \textbf{Binomial},
\textbf{Geometric}, \textbf{Negative Binomial}, \textbf{Hypergeometric}, and
\textbf{Poisson}. The trick to remembering them forever is to never memorize the
formula first. Memorize the \emph{story} first, and let the formula fall out of
the story. By the end you should be able to read a word problem, hear which story
it is telling, and name the distribution before you pick up a pen.

\begin{focusbox}
The single most useful skill in this chapter is \emph{model recognition}: given a
real situation, deciding \emph{which} of these families fits. The questions to ask
are always the same. Is each trial yes/no? Are the trials independent? Is the
number of trials fixed in advance, or are we waiting for something? Are we
sampling with or without replacement? Are we counting rare events in a window of
time or space? Each answer points at a distribution.
\end{focusbox}

\section{Bernoulli: the atom of randomness}
\label{ch07:bernoulli}

Everything in this chapter is built out of one indivisible piece: a single trial
with two outcomes, ``success'' and ``failure.'' Flip a coin. Did a customer buy?
Did the patient recover? Did the bit arrive corrupted? Strip away the story and
you have a number that is $1$ with probability $p$ and $0$ with probability
$1-p$. That is the \textbf{Bernoulli} random variable, named for Jacob Bernoulli,
and it is the atom from which the binomial, geometric, and negative binomial are
all assembled.

\begin{definition}[Bernoulli distribution]
\label{ch07:def-bernoulli}
A random variable $X$ has the \textbf{Bernoulli} distribution with parameter
$p \in [0,1]$, written $X \sim \Bern(p)$, if it takes the value $1$ with
probability $p$ and $0$ with probability $1-p$:
\[
  \P(X = 1) = p, \qquad \P(X = 0) = 1 - p = q.
\]
We call the outcome $X=1$ a \emph{success} and $X=0$ a \emph{failure}; $p$ is the
\emph{success probability} and we often abbreviate $q = 1-p$.
\end{definition}

The choice to code success as the \emph{number} $1$ and failure as $0$ is not
cosmetic. It is exactly what makes the algebra work, because then $X$ is an
\emph{indicator}: $X = \ind_A$ where $A$ is the event ``success.'' Its mean and
variance are immediate.

\begin{proposition}[Mean and variance of Bernoulli]
\label{ch07:bern-moments}
If $X \sim \Bern(p)$ then $\E[X] = p$ and $\Var(X) = p(1-p) = pq$.
\end{proposition}

\begin{proof}
By the definition of expectation for a discrete variable,
\[
  \E[X] = 1 \cdot p + 0 \cdot (1-p) = p.
\]
Since $X$ only takes the values $0$ and $1$, we have $X^2 = X$, so
$\E[X^2] = \E[X] = p$. Therefore
\[
  \Var(X) = \E[X^2] - (\E[X])^2 = p - p^2 = p(1-p) = pq. \qedhere
\]
\end{proof}

\begin{intuitionbox}
The variance $p(1-p)$ is a little parabola, zero at $p=0$ and $p=1$, peaking at
$p = \tfrac12$ where it equals $\tfrac14$. This matches intuition perfectly: a
coin that always lands heads ($p=1$) or always tails ($p=0$) has \emph{no}
randomness, so zero variance. The most unpredictable coin is the fair one, and
that is exactly where the variance is largest.
\end{intuitionbox}

\begin{figure}[H]
\centering
\begin{tikzpicture}
\begin{axis}[
  width=11cm, height=5.2cm,
  xlabel={success probability $p$}, ylabel={$\Var(X) = p(1-p)$},
  xmin=0, xmax=1, ymin=0, ymax=0.28,
  xtick={0,0.25,0.5,0.75,1}, ytick={0,0.1,0.2,0.25},
  title={Bernoulli variance: a parabola peaking at the fair coin},
  samples=100, domain=0:1,
]
\addplot[distblue, very thick] {x*(1-x)};
\addplot[densityred, dashed, thick] coordinates {(0.5,0) (0.5,0.25)};
\node[anchor=south] at (axis cs:0.5,0.25) {\small $\tfrac14$ at $p=\tfrac12$};
\end{axis}
\end{tikzpicture}
\caption{The Bernoulli variance $p(1-p)$ as a function of $p$. It vanishes at the
deterministic extremes $p=0$ and $p=1$ and is largest, equal to $\tfrac14$, at the
fair coin $p=\tfrac12$.}
\label{ch07:fig-bern-var}
\end{figure}

\begin{whybox}
Why does the trick $X^2 = X$ work? Because the only values $X$ can take are $0$
and $1$, and both satisfy $x^2 = x$. Squaring a $0/1$ variable changes nothing.
This is the single most useful identity for indicators, and you will use it
constantly: the variance of an indicator of an event $A$ is
$\P(A)\bigl(1 - \P(A)\bigr)$, full stop.
\end{whybox}

A Bernoulli on its own is almost too simple to be interesting. What makes it
matter is that it is a building block. The next move, which drives the rest of
this chapter, is to line up many independent Bernoulli trials and ask different
questions about them. Each question gives a different named distribution.

\section{Binomial: counting successes in a fixed number of trials}
\label{ch07:binomial}

Here is the first compound story. Fix the number of trials at $n$. Run $n$
independent Bernoulli($p$) trials -- $n$ coin flips, $n$ patients, $n$ shots on
goal. Now ask: \emph{how many successes did I get?} That count, anywhere from $0$
to $n$, is a \textbf{Binomial} random variable.

\begin{intuitionbox}
The keyword is \emph{fixed $n$, count the successes}. You decided in advance how
many trials to run, the trials don't influence each other, each has the same
success chance $p$, and your random quantity is the tally of successes. Whenever a
problem matches that template -- $20$ free throws, $500$ users shown an ad, $8$
births at a hospital -- reach for the binomial.
\end{intuitionbox}

\begin{definition}[Binomial distribution]
\label{ch07:def-binomial}
Let $X$ be the number of successes in $n$ independent $\Bern(p)$ trials. Then $X$
has the \textbf{Binomial} distribution with parameters $n$ and $p$, written
$X \sim \Bin(n,p)$, with PMF
\[
  \P(X = k) = \binom{n}{k} p^{k} (1-p)^{\,n-k}, \qquad k = 0, 1, \dots, n.
\]
\end{definition}

\begin{whybox}
Where does the formula come from? Build it from the story. One \emph{specific}
sequence of $k$ successes and $n-k$ failures -- say SSFSF\dots -- has probability
$p^k (1-p)^{n-k}$, because the trials are independent so we multiply, and there
are $k$ factors of $p$ and $n-k$ factors of $q$. But we don't care \emph{which}
trials succeeded, only \emph{how many}. The number of distinct sequences with
exactly $k$ successes is the number of ways to choose which $k$ of the $n$ slots
are successes, namely $\binom{n}{k}$. Each such sequence has the same probability
$p^k q^{n-k}$, and they are disjoint, so we add them up:
$\binom{n}{k} p^k q^{n-k}$. The two ingredients are: a \emph{counting} factor
$\binom{n}{k}$ (how many arrangements) times a \emph{probability} factor
$p^k q^{n-k}$ (the chance of any one arrangement).
\end{whybox}

Check that these probabilities sum to $1$; the check reveals where the name
``binomial'' comes from. Summing over all $k$,
\[
  \sum_{k=0}^{n} \binom{n}{k} p^{k} q^{\,n-k} = (p + q)^{n} = 1^{n} = 1,
\]
by the binomial theorem. The PMF is literally the binomial expansion of
$(p+q)^n$, term by term. That is the whole reason for the name.

\begin{figure}[H]
\centering
\begin{tikzpicture}
\begin{axis}[
  width=12cm, height=6.2cm,
  ybar, bar width=10pt,
  xlabel={number of successes $k$}, ylabel={$\P(X=k)$},
  xtick={0,1,...,10}, ymin=0, ymax=0.32,
  title={Binomial PMFs for $n=10$: shape shifts with $p$},
  legend style={at={(0.98,0.97)},anchor=north east,font=\small},
  enlarge x limits=0.06,
]
\addplot+[fill=distblue!40,draw=distblue] coordinates {
  (0,0.0282) (1,0.1211) (2,0.2335) (3,0.2668) (4,0.2001) (5,0.1029)
  (6,0.0368) (7,0.0090) (8,0.0014) (9,0.0001) (10,0.0000)};
\addplot+[fill=densityred!35,draw=densityred] coordinates {
  (0,0.0010) (1,0.0098) (2,0.0439) (3,0.1172) (4,0.2051) (5,0.2461)
  (6,0.2051) (7,0.1172) (8,0.0439) (9,0.0098) (10,0.0010)};
\legend{$p=0.3$, $p=0.5$}
\end{axis}
\end{tikzpicture}
\caption{Two binomial PMFs with $n=10$. With $p=0.5$ the bars are symmetric and
centered at $5 = np$. With $p=0.3$ the whole pile shifts left to center near
$3 = np$ and develops a right skew. The expected count $np$ is always where the
mass piles up.}
\label{ch07:fig-binom-pmf}
\end{figure}

\subsection{Mean and variance the easy way}
\label{ch07:binom-moments}

You could compute $\E[X] = \sum_k k \binom{n}{k} p^k q^{n-k}$ directly, and it can
be done, but it is a slog. The elegant route is to remember the story: a binomial
is a \emph{sum} of Bernoullis. Write $X = X_1 + X_2 + \dots + X_n$ where $X_i$ is
the indicator that trial $i$ was a success. Each $X_i \sim \Bern(p)$ and they are
independent.

\begin{proposition}[Mean and variance of Binomial]
\label{ch07:binom-mv}
If $X \sim \Bin(n,p)$ then
\[
  \E[X] = np, \qquad \Var(X) = np(1-p) = npq.
\]
\end{proposition}

\begin{proof}
Write $X = \sum_{i=1}^{n} X_i$ with $X_i \sim \Bern(p)$ independent. Expectation
is linear, \emph{always}, even without independence, so
\[
  \E[X] = \sum_{i=1}^{n} \E[X_i] = \sum_{i=1}^{n} p = np.
\]
For the variance we use that the variance of a sum of \emph{independent} variables
is the sum of the variances:
\[
  \Var(X) = \sum_{i=1}^{n} \Var(X_i) = \sum_{i=1}^{n} p(1-p) = np(1-p). \qedhere
\]
\end{proof}

\begin{happeningbox}
Notice the division of labor in that proof. The mean came from \emph{linearity of
expectation}, which never needs independence -- it would hold even if the trials
were correlated. The variance, on the other hand, leaned on independence: only
because the $X_i$ are independent do the variances simply add with no covariance
terms. This is a recurring theme. Means are robust and add up no matter what;
variances are delicate and add up only when things don't interfere. If your
trials were positively correlated, the count would be \emph{more} spread out than
$npq$; if negatively correlated, less. The hypergeometric, coming up, is exactly a
case of negatively correlated trials, and you will see its variance come out
\emph{smaller} than the binomial's for that reason.
\end{happeningbox}

\begin{example}[Free throws]
\label{ch07:ex-freethrows}
A player makes each free throw independently with probability $p = 0.8$ and takes
$n = 12$ shots. Let $X$ be the number made. Then $X \sim \Bin(12, 0.8)$.
\begin{itemize}
\item Expected made: $\E[X] = np = 12 \times 0.8 = 9.6$.
\item Standard deviation: $\SD(X) = \sqrt{npq} = \sqrt{12 \cdot 0.8 \cdot 0.2}
  = \sqrt{1.92} \approx 1.39$.
\item Probability of making exactly $10$:
  $\binom{12}{10}(0.8)^{10}(0.2)^{2} = 66 \cdot 0.10737 \cdot 0.04
  \approx 0.2835$.
\item Probability of making all $12$: $(0.8)^{12} \approx 0.0687$.
\end{itemize}
So even an $80\%$ shooter sinks all twelve only about $7\%$ of the time -- a nice
gut-check against the intuition that ``a good shooter should hit them all.''
\end{example}

\begin{example}[A/B test conversions]
\label{ch07:ex-abtest}
You show a new checkout button to $n = 500$ visitors. Historically the conversion
rate is $p = 0.06$. If the button doesn't change behavior, the number of
conversions is $X \sim \Bin(500, 0.06)$, with $\E[X] = 30$ and
$\SD(X) = \sqrt{500 \cdot 0.06 \cdot 0.94} \approx 5.31$. If you actually observe
$45$ conversions, that is
$(45 - 30)/5.31 \approx 2.8$ standard deviations above the mean -- a result this
extreme by chance alone is unusual, which is the seed of a hypothesis test. The
binomial is the null-hypothesis model that A/B testing rests on.
\end{example}

\begin{historybox}
The Bernoulli trial and the binomial both trace to Jacob Bernoulli's
\emph{Ars Conjectandi}, published posthumously in 1713. There Bernoulli proved the
first law of large numbers: as you run more and more independent trials, the
observed fraction of successes converges to the true $p$. He saw clearly that the
binomial count $\Bin(n,p)$ concentrates ever more tightly around $np$ as $n$ grows
-- our variance $np(1-p)$ grows only like $n$, while the count itself grows like
$n$, so the \emph{proportion} has variance $p(1-p)/n \to 0$. That single
observation, that averages stabilize, is the foundation on which all of
statistics, polling, and empirical science rests.
\end{historybox}

\begin{trapbox}
The classic binomial mistake is forgetting the counting factor $\binom{n}{k}$.
``The probability of $10$ makes out of $12$ is $0.8^{10} \cdot 0.2^2$'' is wrong;
that is the probability of \emph{one particular} sequence (the first ten in, last
two missed). There are $\binom{12}{10} = 66$ such sequences, and you want any of
them. A second trap: applying the binomial when trials are \emph{not}
independent or $p$ \emph{changes} between trials. Sampling people without
replacement from a small group, for instance, is hypergeometric, not binomial.
\end{trapbox}

\section{Geometric: waiting for the first success}
\label{ch07:geometric}

The binomial fixed $n$ and counted successes. Now flip the question on its head.
Don't fix the number of trials; instead keep going until you get your first
success, and ask: \emph{how many trials did that take?} You are waiting for
something -- the first defective part, the first ``yes,'' the first six on a die.
The waiting time (in trials) is a \textbf{Geometric} random variable.

\begin{intuitionbox}
The keyword is \emph{wait for the first success}. The number of trials is now
random, not fixed; it is the thing you are measuring. Where the binomial answers
``how many successes in $n$ tries,'' the geometric answers ``how many tries until
the first success.''
\end{intuitionbox}

\begin{definition}[Geometric distribution]
\label{ch07:def-geometric}
Run independent $\Bern(p)$ trials with $p > 0$ until the first success. Let $X$ be
the number of trials needed, including the successful one. Then $X$ has the
\textbf{Geometric} distribution with parameter $p$, written $X \sim \Geom(p)$,
with PMF
\[
  \P(X = k) = (1-p)^{\,k-1} p, \qquad k = 1, 2, 3, \dots
\]
\end{definition}

The story makes the formula transparent: to have your first success on trial $k$,
you need the first $k-1$ trials to fail (probability $q^{k-1}$, multiplying
because they're independent) and then the $k$-th to succeed (probability $p$).
Hence $q^{k-1} p$. These sum to $1$ as they should, since it is a geometric series:
\[
  \sum_{k=1}^{\infty} q^{k-1} p = p \sum_{j=0}^{\infty} q^{j}
  = p \cdot \frac{1}{1-q} = \frac{p}{p} = 1.
\]
That geometric series is, of course, where the distribution gets its name.

\begin{marginnotebox}
Beware of two conventions. Some books (and software) define the geometric as the
number of \emph{failures before} the first success, $Y = X - 1$, supported on
$0,1,2,\dots$ with PMF $q^k p$. Both are called ``geometric.'' We use the
\emph{number of trials} version (support starting at $1$) throughout this book, so
its mean is $1/p$, not $q/p$. Always check which convention a source uses.
\end{marginnotebox}

\begin{figure}[H]
\centering
\begin{tikzpicture}
\begin{axis}[
  width=12cm, height=5.6cm,
  ybar, bar width=11pt,
  xlabel={trial of first success $k$}, ylabel={$\P(X=k)$},
  xtick={1,2,...,10}, ymin=0, ymax=0.45,
  title={Geometric PMF, $p=0.4$: a memoryless geometric decay},
  enlarge x limits=0.06,
]
\addplot+[fill=quotientorange!45,draw=quotientorange] coordinates {
  (1,0.4000) (2,0.2400) (3,0.1440) (4,0.0864) (5,0.0518)
  (6,0.0311) (7,0.0187) (8,0.0112) (9,0.0067) (10,0.0040)};
\end{axis}
\end{tikzpicture}
\caption{Geometric PMF with $p=0.4$. The bars decay by a constant factor $q=0.6$
at each step -- each bar is $60\%$ of the one before. The mode is always at
$k=1$: the single most likely outcome is success on the very first try, no matter
how small $p$ is.}
\label{ch07:fig-geom-pmf}
\end{figure}

\subsection{Mean, variance, and the memoryless property}
\label{ch07:geom-mem}

\begin{proposition}[Mean and variance of Geometric]
\label{ch07:geom-mv}
If $X \sim \Geom(p)$ then
\[
  \E[X] = \frac{1}{p}, \qquad \Var(X) = \frac{1-p}{p^{2}}.
\]
\end{proposition}

\begin{proof}
We give the slick first-step argument for the mean; it is more illuminating than
grinding the series. Condition on the first trial. With probability $p$ it
succeeds and we are done in one trial. With probability $q$ it fails, and then --
here is the key -- we are starting over, having already used one trial. The
\emph{remaining} number of trials is again $\Geom(p)$ by independence. So
\[
  \E[X] = p \cdot 1 + q \cdot (1 + \E[X]).
\]
Solve: $\E[X] = p + q + q\E[X]$, hence $\E[X](1-q) = 1$, so $\E[X] = 1/p$. The
variance follows from a similar first-step relation on $\E[X^2]$ (Exercise),
giving $\Var(X) = (1-p)/p^2$.
\end{proof}

\begin{intuitionbox}
The mean $1/p$ is wonderfully intuitive. If a success happens one time in ten
($p = 0.1$), you expect to wait about ten trials to see one. If $p = 1/6$ (rolling
a six), you expect $6$ rolls. The rarer the success, the longer the wait, in exact
reciprocal proportion. This $\E[X] = 1/p$ is one of the most quoted facts in all
of applied probability.
\end{intuitionbox}

The geometric distribution has a famous and almost magical property that no other
discrete waiting-time distribution shares.

\begin{theorem}[Memorylessness]
\label{ch07:geom-memoryless}
If $X \sim \Geom(p)$, then for all integers $s, t \ge 0$,
\[
  \P(X > s + t \given X > s) = \P(X > t).
\]
\end{theorem}

\begin{proof}
First, $\P(X > m) = q^{m}$: to exceed $m$ trials means the first $m$ trials all
failed, probability $q^m$. Then
\[
  \P(X > s+t \given X > s)
  = \frac{\P(X > s+t,\ X > s)}{\P(X > s)}
  = \frac{\P(X > s+t)}{\P(X > s)}
  = \frac{q^{s+t}}{q^{s}} = q^{t} = \P(X > t),
\]
using that $\{X > s+t\} \subseteq \{X > s\}$. \qed
\end{proof}

\begin{paradoxbox}[The gambler's fallacy, formalized]
Memorylessness says: \emph{the dice have no memory}. Suppose you have rolled a die
$20$ times without a six. How long must you now expect to wait? The seductive
wrong answer is ``a six is overdue, it should come soon.'' The correct answer is
that the future is statistically identical to a fresh start: you still expect $6$
more rolls, exactly as if the $20$ misses never happened. The coin, the die, the
roulette wheel -- none of them are keeping a ledger that they owe you a win. The
geometric distribution is the precise mathematical statement that past failures
carry zero information about when the next success arrives. The ``hot hand''
intuition and the ``law of averages'' intuition are both dismantled by this one
equation. (The continuous cousin of this property belongs to the exponential
distribution, which is why exponential and geometric are spiritual siblings -- one
in continuous time, one in discrete trials.)
\end{paradoxbox}

\begin{whybox}
\emph{Why} is the geometric memoryless, intuitively? Because the trials are
independent and identical. The process has no internal state that records what has
happened; each new trial is a clean coin flip with the same $p$. ``How many more
trials until success'' depends only on the same coin you have been flipping all
along, not on how long you have been flipping it. Any waiting-time distribution
built from \emph{independent identical} trials \emph{must} be memoryless -- and in
fact the geometric is the \emph{only} discrete distribution on
$\{1,2,3,\dots\}$ with this property, which is a clean characterization.
\end{whybox}

\begin{example}[Hiring search]
\label{ch07:ex-hiring}
Suppose each candidate you interview is a good hire with probability $p = 0.15$,
independently. Let $X$ be the number of interviews until your first good hire.
Then $X \sim \Geom(0.15)$.
\begin{itemize}
\item Expected interviews: $\E[X] = 1/0.15 \approx 6.67$.
\item Probability the first good hire is the third candidate:
  $\P(X = 3) = (0.85)^2 (0.15) \approx 0.1084$.
\item Probability you need more than $10$ interviews:
  $\P(X > 10) = (0.85)^{10} \approx 0.1969$ -- about a $20\%$ chance of a long
  search.
\end{itemize}
\end{example}

\section{Negative Binomial: waiting for the $r$-th success}
\label{ch07:negbinom}

The geometric waits for the \emph{first} success. The natural generalization
waits for the $r$-th. Keep running independent $\Bern(p)$ trials until you have
accumulated $r$ successes, and ask how many trials that took. That count is a
\textbf{Negative Binomial} random variable. Setting $r = 1$ recovers the
geometric exactly, so the geometric is the special case ``wait for one success.''

\begin{intuitionbox}
Keyword: \emph{wait for the $r$-th success}. Collecting $r$ defects before
shutting down a line; selling $r$ subscriptions; a pitcher recording the $r$-th
strikeout. The number of trials is random, and you stop the instant the $r$-th
success lands.
\end{intuitionbox}

\begin{definition}[Negative Binomial distribution]
\label{ch07:def-negbinom}
Run independent $\Bern(p)$ trials until the $r$-th success ($r \ge 1$ an integer).
Let $X$ be the total number of trials. Then $X$ has the \textbf{Negative
Binomial} distribution with parameters $r$ and $p$, with PMF
\[
  \P(X = k) = \binom{k-1}{r-1} p^{r} (1-p)^{\,k-r}, \qquad k = r, r+1, r+2, \dots
\]
\end{definition}

\begin{whybox}
Build the formula from the story, and watch how it differs from the binomial in
one crucial way. To have the $r$-th success land on \emph{exactly} trial $k$, two
things must happen. First, among the first $k-1$ trials there must be exactly
$r-1$ successes (and $k-r$ failures); second, trial $k$ itself must be a success.
The last trial is \emph{pinned} -- it is forced to be a success, because that's
what stops the process. So we only get to arrange the other $r-1$ successes among
the first $k-1$ slots, giving $\binom{k-1}{r-1}$ ways, not $\binom{k}{r}$. Each
such full sequence has $r$ successes and $k-r$ failures, probability
$p^r q^{k-r}$. Multiply: $\binom{k-1}{r-1} p^r q^{k-r}$. The pinned last trial is
exactly what distinguishes the negative binomial from a binomial; forgetting it
gives the wrong binomial coefficient.
\end{whybox}

Now the easy-way insight pays off again. A negative binomial is a \emph{sum of
$r$ independent geometrics}: wait for the first success ($\Geom(p)$ trials), then
from scratch wait for the second ($\Geom(p)$ more, by memorylessness), and so on
$r$ times. The total trial count is $X = G_1 + G_2 + \dots + G_r$ with each
$G_i \sim \Geom(p)$ independent. Mean and variance fall out instantly by adding.

\begin{proposition}[Mean and variance of Negative Binomial]
\label{ch07:negbinom-mv}
If $X$ is negative binomial with parameters $r$ and $p$, then
\[
  \E[X] = \frac{r}{p}, \qquad \Var(X) = \frac{r(1-p)}{p^{2}}.
\]
\end{proposition}

\begin{proof}
Write $X = \sum_{i=1}^{r} G_i$ with $G_i \sim \Geom(p)$ independent. Then by
linearity $\E[X] = r \cdot (1/p) = r/p$, and by independence
$\Var(X) = r \cdot (1-p)/p^2$. \qed
\end{proof}

\begin{intuitionbox}
If one success takes $1/p$ trials on average, then $r$ successes take $r/p$ -- you
just stack $r$ independent waits end to end. There is nothing to memorize beyond
the geometric; the negative binomial is geometrics in series.
\end{intuitionbox}

\begin{example}[Door-to-door sales]
\label{ch07:ex-sales}
Each house you visit results in a sale with probability $p = 0.2$, independently.
You want $r = 3$ sales. Let $X$ be the number of houses visited.
\begin{itemize}
\item Expected houses: $\E[X] = r/p = 3/0.2 = 15$.
\item Probability the third sale comes exactly at the $10$th house:
  $\P(X = 10) = \binom{9}{2}(0.2)^3 (0.8)^7
  = 36 \cdot 0.008 \cdot 0.2097 \approx 0.0604$.
\item Standard deviation: $\SD(X) = \sqrt{3 \cdot 0.8 / 0.04}
  = \sqrt{60} \approx 7.75$ houses -- a wide spread, so plan for a long day.
\end{itemize}
\end{example}

\begin{figure}[H]
\centering
\begin{tikzpicture}
\begin{axis}[
  width=12cm, height=5.6cm,
  ybar, bar width=8pt,
  xlabel={number of trials $k$}, ylabel={$\P(X=k)$},
  xtick={3,4,...,15}, ymin=0, ymax=0.17,
  title={Negative binomial PMF, $r=3$, $p=0.4$},
  enlarge x limits=0.04,
]
\addplot+[fill=structurepurple!40,draw=structurepurple] coordinates {
  (3,0.0640)(4,0.1152)(5,0.1382)(6,0.1382)(7,0.1244)(8,0.1045)(9,0.0836)
  (10,0.0645)(11,0.0484)(12,0.0355)(13,0.0255)(14,0.0181)(15,0.0127)};
\end{axis}
\end{tikzpicture}
\caption{Negative binomial PMF for $r=3$ successes at $p=0.4$. The support starts
at $k=r=3$ (you cannot finish before the third success), the mass peaks near
$\E[X]=r/p=7.5$, and the long right tail is the signature of a waiting time --
unlucky runs can drag on. Compare to the geometric ($r=1$) of
Figure~\ref{ch07:fig-geom-pmf}, which is this picture with its mode pinned to the
first trial.}
\label{ch07:fig-negbinom-pmf}
\end{figure}

\begin{applicationbox}[Overdispersed counts in data science]
There is a second, statistical reason the negative binomial matters enormously.
When you model count data -- numbers of website visits, insurance claims, reads of
a gene in RNA-sequencing -- the default model is the Poisson (next section), which
forces $\text{variance} = \text{mean}$. Real count data is almost always
\emph{overdispersed}: its variance exceeds its mean. The negative binomial,
re-parameterized to allow any mean and a separately tunable variance, is the
standard fix. In fact a negative binomial arises exactly when the Poisson rate
itself varies randomly (a Gamma-distributed rate), a so-called Gamma--Poisson
mixture. This is why ``negative binomial regression'' is a workhorse in
epidemiology, genomics (the tool DESeq2 is built on it), and actuarial science.
The waiting-time story and the overdispersed-count story are two faces of the same
distribution.
\end{applicationbox}

\section{Hypergeometric: sampling without replacement}
\label{ch07:hypergeometric}

Every distribution so far assumed the trials were \emph{independent} with a
\emph{fixed} success probability $p$. That assumption quietly breaks the moment
you sample from a finite population \emph{without replacement}. Pull a card from a
deck and don't put it back, and the composition of what's left has changed: the
probabilities on the next draw depend on what you already drew. The trials are now
dependent. The count of successes in such a sample follows the
\textbf{Hypergeometric} distribution.

\begin{intuitionbox}
Keyword: \emph{sampling without replacement from a finite pool}. There is a box
with $N$ items, of which $K$ are ``special'' (defective, red, tagged, aces). You
grab a handful of $n$ without replacement. How many specials did you get? That
count is hypergeometric. The tell-tale sign is a \emph{finite population} and
\emph{no replacement} -- if you replaced each item, or if the population were
effectively infinite, you'd be back to the binomial.
\end{intuitionbox}

\begin{definition}[Hypergeometric distribution]
\label{ch07:def-hypergeom}
A population of $N$ items contains $K$ successes and $N-K$ failures. Draw a sample
of $n$ items uniformly at random without replacement. Let $X$ be the number of
successes drawn. Then $X$ has the \textbf{Hypergeometric} distribution with
parameters $N, K, n$, and PMF
\[
  \P(X = k) = \frac{\dbinom{K}{k}\dbinom{N-K}{\,n-k\,}}{\dbinom{N}{n}},
  \qquad \max(0,\, n-(N-K)) \le k \le \min(n, K).
\]
\end{definition}

\begin{whybox}
This is pure counting -- no products of $p$'s at all, because every possible
sample of size $n$ is equally likely. The denominator $\binom{N}{n}$ is the total
number of ways to choose the sample. The numerator counts the favorable samples:
to get exactly $k$ successes, choose which $k$ of the $K$ successes you took
($\binom{K}{k}$ ways) \emph{and} which $n-k$ of the $N-K$ failures you took
($\binom{N-K}{n-k}$ ways), and multiply. Divide favorable by total and you have
the probability. The hypergeometric is the binomial's combinatorial twin: the
binomial multiplies probabilities of independent trials, the hypergeometric
counts equally likely subsets.
\end{whybox}

\begin{proposition}[Mean and variance of Hypergeometric]
\label{ch07:hypergeom-mv}
Let $p = K/N$ be the population success fraction. If $X$ is hypergeometric with
parameters $N, K, n$, then
\[
  \E[X] = n p = n\frac{K}{N}, \qquad
  \Var(X) = n p (1-p) \cdot \frac{N-n}{N-1}.
\]
\end{proposition}

\begin{proof}[Proof sketch]
Write $X = \sum_{i=1}^{n} X_i$ where $X_i$ indicates that the $i$-th draw is a
success. By symmetry, every draw is equally likely to land on any item, so
$\P(X_i = 1) = K/N = p$ for \emph{every} $i$ -- even though the draws are
dependent. Linearity of expectation (which never cares about dependence!) gives
$\E[X] = np$, the same mean as a binomial with the same $p$. The variance is
where dependence bites. The draws are \emph{negatively} correlated: drawing a
success leaves fewer successes for later draws. The covariance terms are negative,
shrinking the variance by exactly the factor $(N-n)/(N-1)$ relative to the
binomial $np(1-p)$. \qed
\end{proof}

\begin{happeningbox}
That extra factor $\dfrac{N-n}{N-1}$ is called the \textbf{finite population
correction}, and it is the whole personality of the hypergeometric. Stare at it:
\begin{itemize}
\item If $n = 1$ (one draw) it equals $1$: a single draw is just a Bernoulli,
  identical to the binomial case. Dependence hasn't had a chance to act.
\item If $n = N$ (sample the entire population) it equals $0$: the variance
  collapses to zero, which is exactly right -- if you draw \emph{everything}, you
  get \emph{all} $K$ successes with certainty, no randomness left.
\item When $N$ is huge compared to $n$, the factor is $\approx 1$ and the
  hypergeometric becomes indistinguishable from $\Bin(n, K/N)$. Removing a few
  fish from a vast lake barely changes the proportions, so with-vs-without
  replacement stops mattering.
\end{itemize}
The correction is the precise mathematical statement of when ``without
replacement'' matters and when you can ignore it.
\end{happeningbox}

\begin{example}[Quality control]
\label{ch07:ex-qc}
A shipment has $N = 50$ chips, of which $K = 5$ are defective. An inspector samples
$n = 10$ without replacement and rejects the shipment if any are defective. Let
$X$ be the number of defectives found, $X \sim$ Hypergeometric$(50, 5, 10)$.
\begin{itemize}
\item Probability of finding zero defectives (and wrongly accepting):
  \[
    \P(X = 0) = \frac{\binom{5}{0}\binom{45}{10}}{\binom{50}{10}}
    = \frac{\binom{45}{10}}{\binom{50}{10}} \approx 0.3106.
  \]
\item Expected defectives found: $\E[X] = 10 \cdot \tfrac{5}{50} = 1$.
\end{itemize}
So there is roughly a $31\%$ chance the sampling plan misses all defectives -- a
sobering number for anyone designing acceptance sampling.
\end{example}

\begin{figure}[H]
\centering
\begin{tikzpicture}
\begin{axis}[
  width=12cm, height=5.8cm,
  ybar, bar width=9pt,
  xlabel={number of successes $k$}, ylabel={$\P(X=k)$},
  xtick={0,1,...,5}, ymin=0, ymax=0.50,
  title={Hypergeometric vs.\ binomial: without vs.\ with replacement},
  legend style={at={(0.98,0.97)},anchor=north east,font=\small},
  enlarge x limits=0.12,
]
\addplot+[fill=targetgreen!40,draw=targetgreen] coordinates {
  (0,0.3106)(1,0.4313)(2,0.2098)(3,0.0442)(4,0.0040)(5,0.0001)};
\addplot+[fill=distblue!35,draw=distblue] coordinates {
  (0,0.3487)(1,0.3874)(2,0.1937)(3,0.0574)(4,0.0112)(5,0.0015)};
\legend{{Hypergeom $(50,5,10)$}, {$\Bin(10,0.1)$}}
\end{axis}
\end{tikzpicture}
\caption{Drawing $n=10$ from a pool of $N=50$ with $K=5$ successes. Both
distributions share the mean $\E[X]=1$, but the hypergeometric (without
replacement) is more concentrated -- its variance carries the finite-population
factor $(N-n)/(N-1)=40/49\approx 0.82$, pulling probability toward the center and
away from the tails relative to the binomial (with replacement).}
\label{ch07:fig-hyper-pmf}
\end{figure}

\begin{example}[Lottery]
\label{ch07:ex-lottery}
A lottery draws $6$ numbers from $1$ to $49$ without replacement; you hold one
ticket with $6$ chosen numbers. The number of your numbers that match is
hypergeometric with $N = 49$, $K = 6$ (your picks are the ``successes''), $n = 6$
(the drawn numbers). The chance of matching all six:
\[
  \P(X = 6) = \frac{\binom{6}{6}\binom{43}{0}}{\binom{49}{6}}
  = \frac{1}{13{,}983{,}816} \approx 7.15 \times 10^{-8}.
\]
One in fourteen million. The hypergeometric is the exact model of every
draw-without-replacement game of chance.
\end{example}

\begin{historybox}
The word ``hypergeometric'' comes not from probability but from the
\emph{hypergeometric series} of the eighteenth and nineteenth centuries (Euler,
Gauss), a vast generalization of the geometric series whose coefficients are
ratios of products much like the ones in this PMF. The distribution inherited the
name because its probabilities are terms of such a series. It is a reminder that
the named distributions did not arrive labeled and gift-wrapped; their names are
historical accidents, dragged in from the analytic machinery mathematicians
happened to be using. Don't read deep meaning into the names -- read it into the
stories.
\end{historybox}

\section{Poisson: counting rare events in a continuum}
\label{ch07:poisson}

We come to the most surprising and arguably the most useful distribution in the
chapter. The Poisson describes the \emph{number of events that occur in a fixed
interval} of time, space, or volume, when the events happen at some average rate
and don't collide with one another: calls arriving at a switchboard per minute,
raindrops hitting a tile per second, mutations along a strand of DNA, typos per
page, customers entering a shop per hour, photons striking a detector. It has a
single parameter $\lambda$, the average number of events expected in the interval.

\begin{definition}[Poisson distribution]
\label{ch07:def-poisson}
A random variable $X$ has the \textbf{Poisson} distribution with rate parameter
$\lambda > 0$, written $X \sim \Pois(\lambda)$, if
\[
  \P(X = k) = e^{-\lambda} \frac{\lambda^{k}}{k!}, \qquad k = 0, 1, 2, \dots
\]
\end{definition}

These sum to $1$ courtesy of the exponential series, which is the source of the
$e^{-\lambda}$ normalizer:
\[
  \sum_{k=0}^{\infty} e^{-\lambda} \frac{\lambda^{k}}{k!}
  = e^{-\lambda} \sum_{k=0}^{\infty} \frac{\lambda^{k}}{k!}
  = e^{-\lambda} \cdot e^{\lambda} = 1.
\]

\begin{figure}[H]
\centering
\begin{tikzpicture}
\begin{axis}[
  width=12cm, height=6cm,
  ybar, bar width=8pt,
  xlabel={number of events $k$}, ylabel={$\P(X=k)$},
  xtick={0,2,4,6,8,10,12,14}, ymin=0, ymax=0.40,
  title={Poisson PMFs: small $\lambda$ is skewed, large $\lambda$ turns bell-like},
  legend style={at={(0.98,0.97)},anchor=north east,font=\small},
  enlarge x limits=0.04,
]
\addplot+[fill=distblue!45,draw=distblue] coordinates {
  (0,0.3679)(1,0.3679)(2,0.1839)(3,0.0613)(4,0.0153)(5,0.0031)(6,0.0005)(7,0.0001)};
\addplot+[fill=densityred!35,draw=densityred] coordinates {
  (0,0.0067)(1,0.0337)(2,0.0842)(3,0.1404)(4,0.1755)(5,0.1755)(6,0.1462)
  (7,0.1044)(8,0.0653)(9,0.0363)(10,0.0181)(11,0.0082)(12,0.0034)(13,0.0013)(14,0.0005)};
\legend{$\lambda=1$, $\lambda=5$}
\end{axis}
\end{tikzpicture}
\caption{Poisson PMFs. For $\lambda=1$ the mass piles near $0$ and is strongly
right-skewed (note $k=0$ and $k=1$ are equally likely). For $\lambda=5$ the
distribution has shifted right, centered near $\lambda=5$, and is starting to look
like the bell curve -- a preview of the normal approximation.}
\label{ch07:fig-pois-pmf}
\end{figure}

\begin{proposition}[Mean and variance of Poisson]
\label{ch07:pois-mv}
If $X \sim \Pois(\lambda)$ then
\[
  \E[X] = \lambda, \qquad \Var(X) = \lambda.
\]
The mean and the variance are equal.
\end{proposition}

\begin{proof}
For the mean, drop the $k=0$ term (it contributes nothing) and shift the index:
\[
  \E[X] = \sum_{k=0}^{\infty} k\, e^{-\lambda}\frac{\lambda^k}{k!}
        = \sum_{k=1}^{\infty} e^{-\lambda}\frac{\lambda^k}{(k-1)!}
        = \lambda \sum_{j=0}^{\infty} e^{-\lambda}\frac{\lambda^{j}}{j!}
        = \lambda \cdot 1 = \lambda,
\]
substituting $j = k-1$ and recognizing the full PMF summing to $1$. For the
variance, the cleanest route computes $\E[X(X-1)]$ by the same index-shift trick:
\[
  \E[X(X-1)] = \sum_{k=2}^{\infty} k(k-1) e^{-\lambda}\frac{\lambda^k}{k!}
  = \lambda^2 \sum_{j=0}^{\infty} e^{-\lambda}\frac{\lambda^{j}}{j!}
  = \lambda^2.
\]
Then $\E[X^2] = \E[X(X-1)] + \E[X] = \lambda^2 + \lambda$, so
$\Var(X) = \E[X^2] - (\E[X])^2 = \lambda^2 + \lambda - \lambda^2 = \lambda$. \qed
\end{proof}

\begin{aestheticbox}
$\E[X] = \Var(X) = \lambda$ is the Poisson's fingerprint. It gives you a free
diagnostic: compute the sample mean and sample variance of any count data, and if
they are wildly different, the data is \emph{not} Poisson. Variance much bigger
than the mean signals overdispersion (reach for the negative binomial); variance
smaller signals underdispersion. This one-number summary -- a distribution whose
spread is locked to its center -- is both the Poisson's great convenience and its
great limitation.
\end{aestheticbox}

\subsection{The law of rare events: Poisson as the limit of the binomial}
\label{ch07:poisson-limit}

Where does this distribution \emph{come from}? The Poisson is not handed down from
on high; it emerges as a limit of the binomial, and that limit is the deepest idea
in the chapter. It explains why so many unrelated phenomena are Poisson.

\begin{intuitionbox}
Imagine counting events in an interval by chopping the interval into a huge number
$n$ of tiny sub-intervals, each so small that at most one event can occur in it.
Each sub-interval is then a Bernoulli trial: an event happens (probability $p$) or
doesn't. The total count is $\Bin(n, p)$. As you slice finer and finer,
$n \to \infty$ while the chance of an event in any one tiny slice
$p \to 0$, but the \emph{expected total} $np = \lambda$ stays fixed at the average
rate. The binomial in this limit \emph{becomes} the Poisson. That is why the
Poisson governs ``many opportunities, each rarely taken.''
\end{intuitionbox}

\begin{theorem}[Poisson limit theorem / law of rare events]
\label{ch07:poisson-limit-thm}
Fix $\lambda > 0$. Suppose $X_n \sim \Bin(n, p_n)$ with $n p_n \to \lambda$ as
$n \to \infty$. Then for each fixed $k = 0,1,2,\dots$,
\[
  \P(X_n = k) \;\longrightarrow\; e^{-\lambda}\frac{\lambda^k}{k!}.
\]
\end{theorem}

\begin{proof}
Set $p_n = \lambda/n$ for clarity (the general case is the same with a vanishing
correction). Then
\[
  \P(X_n = k) = \binom{n}{k}\Bigl(\frac{\lambda}{n}\Bigr)^{k}
  \Bigl(1-\frac{\lambda}{n}\Bigr)^{n-k}.
\]
Expand the binomial coefficient and regroup the powers of $n$:
\[
  \binom{n}{k}\frac{1}{n^k}
  = \frac{n(n-1)\cdots(n-k+1)}{k!\,n^k}
  = \frac{1}{k!}\underbrace{\frac{n}{n}\frac{n-1}{n}\cdots\frac{n-k+1}{n}}_{\to\,1}.
\]
For fixed $k$, each of the $k$ fractions tends to $1$ as $n\to\infty$, so this
whole piece tends to $1/k!$. Meanwhile
$\bigl(1-\lambda/n\bigr)^{n} \to e^{-\lambda}$ (the defining limit of the
exponential), and $\bigl(1-\lambda/n\bigr)^{-k} \to 1$. Collecting the surviving
pieces,
\[
  \P(X_n = k) \;\longrightarrow\; \frac{1}{k!}\cdot \lambda^k \cdot e^{-\lambda}
  = e^{-\lambda}\frac{\lambda^k}{k!}. \qed
\]
\end{proof}

\begin{historybox}
Sim\'eon Denis Poisson published this distribution in 1837 in a treatise on the
probability of judicial verdicts, but it languished as a mathematical curiosity
for decades. Its fame came in 1898 when Ladislaus Bortkiewicz studied a famously
morbid dataset: the number of Prussian cavalry soldiers killed each year by being
kicked by a horse, recorded across $14$ corps over $20$ years. Deaths-by-horse are
the perfect Poisson scenario: enormously many soldiers ($n$ huge), each with a
tiny chance of this particular fate ($p$ tiny), giving a modest expected count
($\lambda \approx 0.7$ per corps-year). Bortkiewicz showed the yearly counts
matched the Poisson almost exactly, and his little book \emph{The Law of Small
Numbers} gave the law of rare events its name. The deep point: \emph{the
mechanism, not the subject matter, makes something Poisson}. Horse-kicks, radio
decays, and typos obey the same law because they share the same structure.
\end{historybox}

\begin{example}[Poisson approximation in action]
\label{ch07:ex-poisapprox}
A factory produces widgets with a $0.1\%$ defect rate. In a batch of $n = 2000$,
let $X$ be the number of defective widgets. Exactly, $X \sim \Bin(2000, 0.001)$,
which is painful to compute by hand. But $np = \lambda = 2$, and $p$ is tiny, so
$X \approx \Pois(2)$. The probability of zero defectives:
\[
  \P(X = 0) \approx e^{-2} \approx 0.1353.
\]
The exact binomial value is $(0.999)^{2000} \approx 0.1352$ -- agreement to four
decimals. The probability of at most one defective:
$e^{-2}(1 + 2) = 3e^{-2} \approx 0.406$. The Poisson turns a brutal binomial into
a one-line calculation, and the law of rare events guarantees the answer is
excellent whenever $n$ is large and $p$ small.
\end{example}

\begin{example}[Call center staffing]
\label{ch07:ex-callcenter}
Calls arrive at a help line at an average rate of $\lambda = 4$ per minute. Model
the count in a given minute as $\Pois(4)$.
\begin{itemize}
\item Exactly $4$ calls: $\P(X=4) = e^{-4}\dfrac{4^4}{4!}
  = e^{-4}\dfrac{256}{24} \approx 0.1954$.
\item No calls in a minute: $\P(X = 0) = e^{-4} \approx 0.0183$.
\item More than $6$ calls: $\P(X > 6) = 1 - \sum_{k=0}^{6} e^{-4}\frac{4^k}{k!}
  \approx 1 - 0.889 = 0.111$. Roughly an $11\%$ chance any given minute is a
  surge -- the kind of number that drives staffing decisions.
\end{itemize}
\end{example}

\begin{connectionbox}
The Poisson does not live alone. It is the count side of the \textbf{Poisson
process}, the canonical model of random arrivals in continuous time, which we
build carefully in a later chapter. In that process, the \emph{number} of arrivals
in an interval is Poisson, while the \emph{waiting time between} consecutive
arrivals is Exponential -- the continuous, memoryless cousin of the geometric. So
the discrete waiting-time/count pair (geometric, binomial) has a continuous mirror
image (exponential, Poisson), and the bridge between them is exactly the limiting
argument of this section. When you later meet the exponential and gamma
distributions, remember that they are the continuous-time reincarnations of the
geometric and negative binomial you learned here. The whole landscape rhymes.
\end{connectionbox}

\begin{applicationbox}[Poisson in modern data science]
Poisson assumptions are everywhere in applied work. \emph{Poisson regression}
models count outcomes (hospital admissions, goals scored, clicks) as a function of
predictors, and is the count-data analogue of linear regression. In natural
language processing, simple word-count models lean on Poisson assumptions. In
genomics, read counts in sequencing are Poisson under idealized sampling (and
negative binomial once biological variation is added). In sports analytics, goals
in soccer are famously well modeled as Poisson, which lets you compute the
probability of any score line from two rate parameters. Whenever you find yourself
counting independent rare events in a window, the Poisson is your first model and
your null hypothesis.
\end{applicationbox}

\section{Putting it together: which distribution fits?}
\label{ch07:decision}

The whole point of naming these families is to recognize them in the wild. Here is
the decision tree, distilled. Run a problem through these questions in order.

\begin{focusbox}
\textbf{1. Single yes/no trial?} $\Rightarrow$ \textbf{Bernoulli}$(p)$.

\textbf{2. Fixed number $n$ of independent trials, counting successes?}
$\Rightarrow$ \textbf{Binomial}$(n, p)$.

\textbf{3. Independent trials, waiting for the \emph{first} success (count the
trials)?} $\Rightarrow$ \textbf{Geometric}$(p)$.

\textbf{4. Independent trials, waiting for the \emph{$r$-th} success?}
$\Rightarrow$ \textbf{Negative Binomial}$(r, p)$.

\textbf{5. Drawing a sample \emph{without replacement} from a finite pool,
counting successes?} $\Rightarrow$ \textbf{Hypergeometric}$(N, K, n)$.

\textbf{6. Counting rare events in an interval of time/space at a fixed average
rate?} $\Rightarrow$ \textbf{Poisson}$(\lambda)$.
\end{focusbox}

The two perennial points of confusion are worth isolating.

\begin{trapbox}
\textbf{Binomial vs.\ Hypergeometric.} Both count successes. The difference is
\emph{replacement}. With replacement (or an effectively infinite population) the
trials are independent and $p$ is constant: binomial. Without replacement from a
finite pool, the trials are dependent and $p$ drifts as you draw: hypergeometric.
When the population is much larger than the sample ($N \gg n$), the two coincide,
and the binomial is the easy approximation.

\textbf{Binomial vs.\ Negative Binomial.} Both involve trials and successes. The
difference is \emph{what is fixed and what is random}. Binomial: $n$ is fixed, the
number of successes is random. Negative binomial: the number of successes $r$ is
fixed, and the number of trials is random. ``I ran $20$ trials, how many wins?'' is
binomial. ``I kept going until $5$ wins, how many trials?'' is negative binomial.
\end{trapbox}

\begin{walkbox}
A test has $30$ multiple-choice questions, each with $4$ options, and a student
guesses every one. What is the distribution of the number correct, and what is its
expected value? Pause and run it through the tree before reading on.

\emph{Answer:} fixed $n = 30$ independent trials, each a success (correct guess)
with probability $p = 1/4$, counting successes. That is question 2 on the tree:
$\Bin(30, 1/4)$, with $\E[X] = 30 \cdot 1/4 = 7.5$. A pure guesser averages $7.5$
out of $30$.
\end{walkbox}

\begin{table}[H]
\centering
\renewcommand{\arraystretch}{1.5}
\begin{tabular}{@{}lccc@{}}
\toprule
\textbf{Distribution} & \textbf{PMF} & \textbf{Mean} & \textbf{Variance} \\
\midrule
$\Bern(p)$ & $p^k(1-p)^{1-k}$ & $p$ & $p(1-p)$ \\
$\Bin(n,p)$ & $\binom{n}{k}p^k(1-p)^{n-k}$ & $np$ & $np(1-p)$ \\
$\Geom(p)$ & $(1-p)^{k-1}p$ & $\dfrac{1}{p}$ & $\dfrac{1-p}{p^2}$ \\
NegBin$(r,p)$ & $\binom{k-1}{r-1}p^r(1-p)^{k-r}$ & $\dfrac{r}{p}$ & $\dfrac{r(1-p)}{p^2}$ \\
Hypergeom$(N,K,n)$ & $\dfrac{\binom{K}{k}\binom{N-K}{n-k}}{\binom{N}{n}}$ & $n\dfrac{K}{N}$ & $n\dfrac{K}{N}\dfrac{N-K}{N}\dfrac{N-n}{N-1}$ \\
$\Pois(\lambda)$ & $e^{-\lambda}\dfrac{\lambda^k}{k!}$ & $\lambda$ & $\lambda$ \\
\bottomrule
\end{tabular}
\caption{The six discrete families at a glance. Notice the family resemblances:
the binomial mean $np$ and the hypergeometric mean $n(K/N)$ are the same with
$p = K/N$, differing only by the finite-population factor in the variance; the
negative binomial mean $r/p$ is just $r$ copies of the geometric mean $1/p$.}
\label{ch07:tab-summary}
\end{table}

\begin{happeningbox}
Step back and admire the architecture. The Bernoulli is the atom. Sum $n$ of them
and ask ``how many successes'' $\to$ binomial. Ask instead ``how long until the
first success'' $\to$ geometric; ``until the $r$-th'' $\to$ negative binomial.
Break the independence by sampling without replacement $\to$ hypergeometric. Push
the binomial to its rare-event limit $\to$ Poisson. Every distribution in this
chapter is a different \emph{question} asked of the same underlying coin-flipping
process, or a different \emph{limit} of it. You are not memorizing six unrelated
formulas; you are reading six chapters of one story. That is the sense in which you
can now \emph{see} discrete probability rather than recite it.
\end{happeningbox}

\section{Exercises}
\label{ch07:exercises}

\begin{exercise}[Bernoulli warm-up]
\label{ch07:exr-bern}
A biased coin has $\P(\text{heads}) = 0.7$. Let $X = 1$ if heads. Compute $\E[X]$,
$\Var(X)$, and $\SD(X)$. For what value of $p$ is the variance of a Bernoulli
maximized, and what is that maximum?
\end{exercise}

\begin{exercise}[Binomial computations]
\label{ch07:exr-binom}
A multiple-choice quiz has $8$ questions, each with $5$ choices; a student guesses
randomly. Let $X$ be the number correct.
(a) Name the distribution of $X$ with its parameters.
(b) Compute $\P(X = 2)$.
(c) Compute $\E[X]$ and $\Var(X)$.
(d) Compute $\P(X \ge 1)$, the chance of getting at least one right.
\end{exercise}

\begin{exercise}[Geometric waiting]
\label{ch07:exr-geom}
You roll a fair die repeatedly until the first six.
(a) What is the probability the first six appears on the fourth roll?
(b) What is the expected number of rolls?
(c) Given that you have rolled five times with no six, what is the expected number
of \emph{additional} rolls until the first six? Explain how memorylessness makes
this trivial.
\end{exercise}

\begin{exercise}[Negative binomial]
\label{ch07:exr-negbinom}
A basketball player makes each free throw with probability $0.6$. She shoots until
she has made $4$ baskets. Let $X$ be the number of shots taken.
(a) Identify the distribution and its parameters.
(b) Compute $\P(X = 6)$.
(c) Compute $\E[X]$ and $\SD(X)$.
(d) Explain, using the ``sum of geometrics'' picture, why $\E[X] = 4/0.6$.
\end{exercise}

\begin{exercise}[Hypergeometric]
\label{ch07:exr-hypergeom}
A standard $52$-card deck contains $4$ aces. You are dealt a $5$-card poker hand.
Let $X$ be the number of aces in your hand.
(a) Identify the distribution and its parameters.
(b) Compute $\P(X = 2)$.
(c) Compute $\E[X]$.
(d) If instead you drew $5$ cards \emph{with} replacement (reshuffling each time),
what distribution would $X$ follow, and how would $\E[X]$ compare?
\end{exercise}

\begin{exercise}[Poisson modeling]
\label{ch07:exr-pois}
A website receives an average of $3$ sign-ups per hour, modeled as Poisson.
(a) Probability of exactly $5$ sign-ups in a given hour.
(b) Probability of no sign-ups in a given hour.
(c) Probability of at least $2$ sign-ups in a given hour.
(d) What is the distribution of sign-ups in a \emph{half}-hour, and its mean?
(Hint: the rate scales with the length of the interval.)
\end{exercise}

\begin{exercise}[Poisson approximation]
\label{ch07:exr-poisapprox}
Each of $10{,}000$ independent visitors clicks a rare banner ad with probability
$0.0003$.
(a) Give the exact distribution of the number of clicks $X$.
(b) Justify approximating $X$ by a Poisson and give $\lambda$.
(c) Use the Poisson approximation to compute $\P(X = 0)$, $\P(X = 3)$, and
$\P(X \ge 1)$.
\end{exercise}

\begin{exercise}[Conceptual: which model?]
\label{ch07:exr-which}
For each scenario, name the distribution and its parameters, and justify in one
sentence by referencing the decision tree.
(a) The number of defective bulbs in a random sample of $20$ drawn without
replacement from a crate of $200$ containing $15$ defectives.
(b) The number of emails you receive between $9$ and $10$ a.m., averaging $12$.
(c) The number of coin flips until you see two heads in total.
(d) The number of heads in $50$ flips of a fair coin.
(e) Whether a single randomly chosen voter supports a measure.
\end{exercise}

\begin{exercise}[Proof: geometric variance]
\label{ch07:exr-geomvar}
Let $X \sim \Geom(p)$. Using a first-step argument analogous to the one for the
mean, derive $\E[X^2]$ and hence show $\Var(X) = (1-p)/p^2$.
(Hint: condition on the first trial. If it fails, the remaining count $X'$
satisfies $X = 1 + X'$ with $X' \sim \Geom(p)$ independent, so
$\E[X^2] = p\cdot 1 + q\,\E[(1+X')^2]$. Expand and solve, using $\E[X] = 1/p$.)
\end{exercise}

\begin{exercise}[Proof: binomial mode]
\label{ch07:exr-binommode}
Let $X \sim \Bin(n,p)$. By examining the ratio
$\dfrac{\P(X=k)}{\P(X=k-1)}$, show that this ratio exceeds $1$ exactly when
$k < (n+1)p$. Conclude that the PMF increases then decreases, and that the most
likely value (the mode) is $\floor{(n+1)p}$ (with a tie when $(n+1)p$ is an
integer). Verify for the free-throw example $\Bin(12, 0.8)$ that the mode is $10$.
\end{exercise}

\begin{exercise}[Proof: Poisson sum]
\label{ch07:exr-poissum}
Let $X \sim \Pois(\lambda)$ and $Y \sim \Pois(\mu)$ be independent. Prove that
$X + Y \sim \Pois(\lambda + \mu)$.
(Hint: $\P(X+Y = n) = \sum_{k=0}^{n}\P(X=k)\P(Y=n-k)$. Factor out
$e^{-(\lambda+\mu)}/n!$ and recognize the binomial theorem
$\sum_k \binom{n}{k}\lambda^k\mu^{n-k} = (\lambda+\mu)^n$ inside the sum.)
Interpret the result through the law of rare events: merging two independent
streams of rare events gives one rare-event stream at the combined rate.
\end{exercise}

\begin{exercise}[Connecting the families]
\label{ch07:exr-connect}
(a) Show that the negative binomial with $r=1$ is exactly the geometric.
(b) Explain in words why the hypergeometric variance is always $\le$ the binomial
variance with matching $p = K/N$, and identify the case where they are equal.
(c) Two friends each model the number of typos on a page. One uses $\Bin(n,p)$
with $n$ characters and tiny per-character error rate $p$; the other uses
$\Pois(\lambda)$. Under what relationship between the parameters do they
essentially agree, and which theorem of this chapter justifies it?
\end{exercise}
